\pdfbookmark[0]{Notices. Investigations into Topics in Higher Geodesy. 1846.}{anzeigen-geodaesie-zweite}
\begin{center}
\rule{.85\textwidth}{.3pt}\par
Nachrichten von der K. Gesellschaft der Wissenschaften zu Göttingen.\par 28 September 1846.\par
\rule{.85\textwidth}{.3pt}
\end{center}
\hspace*{\parindent}On 1\textsuperscript{st} September, Privy Court Councillor \textsc{Gauss} submitted a lecture to the Royal Society of Sciences under the title
\begin{center}\textit{Investigations into Topics in Higher Geodesy, Second Memoir},\end{center}
\noindent whose contents and connection with the first memoir are to be briefly described here.

The first memoir presented a new method for treating geodetic measurements. Its main feature is that most calculations are carried out just as, or almost as, they would be if the triangulation network lay on a sphere instead of a spheroid, with no loss whatsoever in the utmost accuracy of the results. One of the principal problems in geodetic calculation starts from the length of a triangle side that is a geodesic line, the latitude of one endpoint, and the azimuth at which the other endpoint appears from it. The aim is to determine the latitude of that other endpoint, the azimuth of the side there, and the difference in longitude between the two points. With the method just described, this reduces to solving a spherical triangle. Nevertheless, a couple of pages in that memoir are devoted to the problem because the ordinary formulas of spherical trigonometry would not always give the results quite sufficient accuracy unless logarithms with more digits were used. Certain transformations of those formulas therefore become necessary. Moreover, the sides of triangles whose angles are actually measured can only ever be small quantities compared with the dimensions of the %
